%% -*- Mode: LaTeX -*-
%%
%% svn-to-git.tex
%% Crated Fri May 17 14:27:54 AKDT 2019
%% by Raymond E. Marcil <marcilr@gmail.com>
%% 
%% Subversion to Git Conversion
%%
%% Links
%% =====
%% LaTeX/Modular Documents
%% 'Working on discrete parts of your documents has
%% consequences for how the base file is compiled;
%% these will be dealt with below.'
%% https://en.wikibooks.org/wiki/LaTeX/Modular_Documents
%%

  %%
%%%%%% Preamble.
  %%

%% Specify DVIPS driver used by things like hyperref
\documentclass[12pt,letterpaper,dvips]{article}

%% Set title, author, and date at top of document
\title{Subversion to Git Conversion}
\author{Raymond E. Marcil\\
        \texttt{$<$marcilr@gmail.com$>$}
}

%%
%% Display current date from datetime package
%% http://ctan.math.washington.edu/tex-archive/obsolete/\
%% macros/latex/contrib/datetime/datetime.pdf
%%
\date{\today}


%% rcs is the package to display cvs revision info.
%%\usepackage{rcs}
\usepackage{fullpage}
\usepackage{fancyvrb} 
\usepackage{graphicx}
\usepackage{figsize}
\usepackage{calc}

%%
%% enumitem – Control layout of itemize, enumerate, description
%% https://www.ctan.org/pkg/enumitem
%%
%% Allows for use of \bgein{itemize}[leftmargin=0pt] 
%% to lists with 0 left margin.
%%
%% Itemize left margin
%% http://tex.stackexchange.com/questions/170525/itemize-left-margin
%% 
\usepackage{enumitem}%     http://ctan.org/pkg/enumitem


%% caption package for use in justifying table or figure captions
\usepackage{caption}

%%
%% datetime.sty v2.60: Formatting
%% Current Date and Time
%% Nicola L. C. Talbot
%% Dickimaw Books
%% http://www.dickimaw-books.com/
%% 2015-03-20
%% http://ctan.math.washington.edu/tex-archive/obsolete/\
%% macros/latex/contrib/datetime/datetime.pdf
%%
\usepackage{datetime}

\usepackage{lipsum}
\usepackage{xspace}
\usepackage{booktabs}
\usepackage[first,bottomafter]{draftcopy}
\usepackage[numbib]{tocbibind}

\usepackage{amssymb}              %% AMS Symbols, used for \checkmark
\usepackage{multicol}

%%
%% Hyperref package for embedding URLs for clickable links in PDFs, 
%% also specify PDF attributes here.
%%
%% The pdfborder={0 0 0} is what ellimated the blue box around the url
%% displayed by \href{}{}.
%%
%% The command pdfborder={0 0 1} would display a box with thickness of 1 pt.
%%
%% Hypertext marks in LATEX: a manual for hyperref
%% by Sebastian Rahtz and Heiko Oberdiek - November 2012
%% http://ctan.org/pkg/hyperref 
%% http://mirror.hmc.edu/ctan/macros/latex/contrib/hyperref/doc/manual.html
%%
\usepackage[
colorlinks,
linkcolor=blue,
%%colorlinks=false,
hyperindex=false,
urlcolor=blue,
pdfborder={0 0 0},
pdfauthor={Raymond E. Marcil},
pdftitle={Subversion to Git Conversion},
pdfcreator={ps2pdf},
pdfsubject={LaTeX Examples},
pdfkeywords={LaTeX Examples}
]{hyperref}


  %%
%%%%%% Customization.
  %%

% On letter paper with 10pt font the Verbatim environment has 65 columns.
% With 12pt font the environment has 62 columns.  Exceeding this will exceed
% the frame and will look ugly.  YHBW.  HAND.
\RecustomVerbatimEnvironment{Verbatim}{Verbatim}{frame=single}

\renewenvironment{description}
                 {\list{}{\labelwidth 0pt \iteminden-\leftmargin
                          \let\labelsep\hsize
                          \let\makelabel\descriptionlabel}}
                 {\endlist}
\renewcommand*\descriptionlabel[1]{\hspace\labelsep\sffamily\bfseries #1}


  %%
%%%%%% Commands.
  %%

\newcommand{\FIXME}[1]{\textsf{[FIXME: #1]}}
\newcommand{\cmd}[1]{\texttt{#1}}


%% Squeeze space above/below captions
\setlength{\abovecaptionskip}{4pt}   % 0.5cm as an example
\setlength{\belowcaptionskip}{4pt}   % 0.5cm as an example


%% Tex really adds a lot of whitespace to itemized 
%% lists so define a new command itemize* with a 
%% lot less whitespace.  Found this in the British
%% Tex faq.
\newenvironment{itemize*}%
  {\begin{itemize}%
    \setlength{\itemsep}{0pt}%
    \setlength{\parsep}{0pt}}%
  {\end{itemize}}

  
%%
%% Tex really adds a lot of whitespace to itemized 
%% lists so define a new command itemize* with a 
%% lot less whitespace.  Found this in the British
%% Tex faq.
%%
%% Tue Jun 23 13:22:04 AKDT 2015
%% =============================
%% Added [leftmargin=0.0mm] to set the left margin=0
%% This requires use of the enumitem package:
%%   \usepackage{enumitem}%     http://ctan.org/pkg/enumitem
%%
%% Itemize left margin
%% http://tex.stackexchange.com/questions/170525/itemize-left-margin
%%
\newenvironment{itemizenoleft*}%
  {\begin{itemize}[leftmargin=15.0pt]%
    \setlength{\itemsep}{0pt}%
    \setlength{\parsep}{0pt}}%
  {\end{itemize}}
  

%%
%% Tex really adds a lot of whitespace to itemized 
%% lists so define a new command enumerate* with a 
%% lot less whitespace.  Created using itemize*
%% pattern.  
%%
  \newenvironment{enumerate*}%
  {\begin{enumerate}%
    \setlength{\itemsep}{0pt}%
    \setlength{\parsep}{0pt}}%
  {\end{enumerate}}


%%
%% Tex really adds a lot of whitespace to itemized 
%% lists so define a new command enumerate* with a 
%% lot less whitespace.  Created using itemize*
%% pattern.  
%%
%% Tue Jun 23 13:22:04 AKDT 2015
%% =============================
%% Added [leftmargin=0.0mm] to set the left margin=0
%% This requires use of the enumitem package:
%%   \usepackage{enumitem}%     http://ctan.org/pkg/enumitem
%%
%% Itemize left margin
%% http://tex.stackexchange.com/questions/170525/itemize-left-margin
%%
\newenvironment{enumeratenoleft*}%
  {\begin{enumerate}[leftmargin=0.0mm]%
    \setlength{\itemsep}{0pt}%
    \setlength{\parsep}{0pt}}%
  {\end{enumerate}}


%% Squeeze space
\renewcommand\floatpagefraction{.9}
\renewcommand\topfraction{.9}
\renewcommand\bottomfraction{.9}
\renewcommand\textfraction{.1}   
\setcounter{totalnumber}{50}
\setcounter{topnumber}{50}
\setcounter{bottomnumber}{50}


  %%
%%%%%% Document.
  %%

\begin{document}

%%
%% \input{filename} imports the commands from filename.tex
%% into the target file; it's equivalent to typing all the
%% commands from filename.tex right into the current file
%% where the \input line is.
%%
%% \include{filename} essentially does a \clearpage before
%% and after \input{filename}, together with some magic to
%% switch to another .aux file, and omits the inclusion at
%% all if you have an \includeonly without the filename in
%% the argument.  This is primarily useful when you have a
%% big project on a slow computer; changing one of the
%% include targets won't force you to regenerate the
%% outputs of all the rest.
%%
%% \include{filename} gets you the speed bonus, but it
%% also can't be nested, can't appear in the preamble,
%% and forces page breaks around the included text.
%%
%% --When should I use \input vs. \include?
%%   https://tex.stackexchange.com/questions/246/\
%%   when-should-i-use-input-vs-include
%%
%% NOTE: \newpage is embedded at the beginning of
%%       included sections as needed.
%%
%% -*- Mode: LaTeX -*-
%%
%%
%% titlepage.tex
%% Created Fri May 31 09:55:40 AKDT 2019
%  Copyright (C) 2019 by Raymond E. Marcil <marcilr@gmail.com>
%%

%% ==============================================================
%% Abstract
%% ==============================================================
\maketitle
\begin{abstract}
Details for Subversion to Git revision control system conversion.
\lipsum[1]
\end{abstract}

%% Draw DNR logo and address at bottom of page
\begin{figure}[h]
        \hspace{0.32in}
        \SetFigLayout{1}{1}
        \begin{minipage}[b]{0.16\figwidth}
                \includegraphics[width=\textwidth]{images/dnr_bwlogo.eps}
        \end{minipage}
        \hspace{5pt}
        \begin{minipage}[b]{\figwidth}
                \bf{Alaska Department of Natural Resources}\\
                \small{\sf{Division of Support Services\\
                Land Records Information Section\\
                550 W. 7th Ave. Suite 706\\
                Anchorage, Alaska 99501}}
        \end{minipage}
\end{figure}

%% -*- Mode: LaTeX -*-
%%
%% contents.tex
%% Created Fri May 31 10:28:26 AKDT 2019
%% Copyright (C) 2019 by Raymond E. Marcil <marcilr@gmail.com>
%%
\newpage
\tableofcontents

%% -*- Mode: LaTeX -*-
%%
%% figures.tex
%% Created Fri May 31 10:52:06 AKDT 2019
%% Copyright (C) 2019 by Raymond E. Marcil <marcilr@gmail.com>
%%
%% List of figures
%%
\newpage
\listoffigures

%% -*- Mode: LaTeX -*-
%%
%% tables.tex
%% Created Fri May 31 10:56:50 AKDT 2019
%% Copyright (C) 2019 by Raymond E. Marcil <marcilr@gmail.com>
%%
%% List of tables
%%
\listoftables

%% -*- Mode: LaTeX -*-
%%
%% abbreviations.txt
%% Created Fri May 31 10:46:19 AKDT 2019
%% Copyright (C) 2019 by Raymond E. Marcil <marcilr@gmail.com>
%%

%% =============== List of Abbreviations ===============
%% =============== List of Abbreviations ===============
\newpage
\setcounter{secnumdepth}{0}
\section{List of Definitions and Abbreviations}
\begin{itemize*}
  \item{\begin{bf}MOA\end{bf}} - Municipality of Anchorage

\end{itemize*}

%% -*- Mode: LaTeX -*-
%%
%% introduction.tex
%% Created Fri May 31 11:01:45 AKDT 2019
%% Copyright (C) 2019 by Raymond E. Marcil <marcilr@gmail.com>
%%
%% Introduction
%%

%% ====================== Introduction ===========================
%% ====================== Introduction ===========================
\newpage
\section{Introduction}
Introduction to the \LaTeX\ Template with Examples.

%% -*- Mode: LaTeX -*-
%%
%% examples.tex
%% Created Fri May 31 12:37:05 AKDT 2019
%% Copyright (C) 2019 by Raymond E. Marcil <marcilr@gmail.com>
%%
%% Examples
%%

%% ======================== Examples =============================
%% ======================== Examples =============================
\newpage
\section{Examples}
Series of useful \LaTeX\ markup. Need to break out to 
separate examples.tex file.

\subsection{Escaping $<$ and $>$ Symbols}
To get \$$<$\$ or \$$>$\$ just wrap the symbols in \$ for math mode.

\subsection{Enumerate}
\begin{enumerate}
  \item{DNR} - Alaska State Department of Natural Resources
    \begin{itemize*}
      \item{HI} - Historical Index, not maintained since 1982
      \item{LE} - Land Estate, maintained by SGU
      \item{ME} - Mineral Estate, maintaind by SGU
    \end{itemize*}

  \item{Alaska State Surveys}
    \begin{itemize*}
      \item{ASBLT} - As-Built Survey
      \item{ASCS} - Cadastral Survey
    \end{itemize*}
\end{enumerate}

%%
%%
%% comments.tex
%% Created Fri May 31 12:41:00 AKDT 2019
%% Copyright (C) 2019 by Raymond E. Marcil <marcilr@gmail.com>
%%
%% Comments
%%

%% ======================= Comments =======================
%% ======================= Comments =======================
\subsection{Comments}
\begin{center}
\framebox{
\begin{minipage}[t]{0.9\textwidth}
\cmd{COMMENTS} Comment --- \emph{Sean Weems, Spring 2003}\\
We should get the \cmd{COMMENTS} column searchable via the 
landrecords application before we do much anything else -- shouldn't
be too hard.
\end{minipage}
}
\end{center}

\begin{center}
\framebox{
\begin{minipage}[t]{0.9\textwidth}
\emph{Errata: Plats spanning multiple sections}\\
A few anomalies can be observed in the \cmd{AKPLATS}
table. Specifically plats exist that span multiple sections. 
Since the table only has a single column, \cmd{SCODE}, 
that accepts a single section code, SGU (Status
Graphics Unit) has handled this problem by entering multiple 
rows in the table, each with a different section that point to
the same plat or file. Multiple section plats are indicated by  setting 
the \cmd{TCODE} column to the value \cmd{37}, and making an 
appropriate notation like \emph{Section 24-25-26-27} in the 
\cmd{REMARKS} column.\\
\FIXME{Perhaps the \cmd{SCODE} column should accept an array of sections?}

\end{minipage}
}
\end{center}


%%%% -*- Mode: LaTeX -*-
%%
%% hyperlinks.tex
%% Created Fri May 31 12:48:03 AKDT 2019
%% Copyright (C) 2019 by Raymond E. Marcil <marcilr@gmail.com>
%%

%% ====================== Hyperlinks ======================
%% ====================== Hyperlinks ======================
\subsection{Hyperlinks}
Use $\backslash$href\{\} to generate hyperlinks:

\begin{verbatim}
  \href{http://www.google.com}{Google}}
\end{verbatim}

\noindent Yields: \href{http://www.google.com}{Google}

%%%% -*- Mode: LaTeX -*-
%%
%% table-examples.tex
%% Created Fri May 31 12:52:31 AKDT 2019
%% Copyright (C) 2019 by Raymond E. Marcil <marcilr@gmail.com>
%%
%% Table Examples
%%

%% ================ Table Examples ================
%% ================ Table Examples ================
\newpage
\subsection{Table Examples}

\begin{table}[htb]
\begin{tabular}{|p{.25\textwidth}|p{.20\textwidth}|p{.47\textwidth}|}\hline 
Column Name&Type&Description\\ 
\hline
EQS&VARCHAR2(1)&!NULL map shows village selections\\
ITM\_COL&VARCHAR2(1)&USGS ITM column: 1-6\\
ITM\_ROW&VARCHAR2(1)&USGS ITM row: A-E\\
QMQ\_ABBR\_DNR&VARCHAR2(3)&Three character DNR abbreviation for the QMQ\\
RASTER\_FILENAME&VARCHAR2(50)&Physical path to file\\
RASTER\_PATHNAME&VARCHAR2(50)&URL path to PDF of map\\
SCODE&VARCHAR2(2)&Supplement map code: 1,2,3,...\\
COMMENTS&VARCHAR2(256)&Plat comments\\
\hline
\end{tabular}
\caption {\cmd{EASEMENTS\_17B} Table}
\label{table:easements_17b}
\end{table}

%%
%% This table example was breaking the build.
%% Commnented out for analysis.
%%
%% A simple table example
%% The htb attribute attempts to inline the table or figure
%% where you put it in the document
%%\begin{table}[htb]
%%\begin{center}
%%\begin{tabular*}{\textwidth}{@{}p{.25\textwidth}@{}p{.75\textwidth}}
%%\hline
%%\hline\\[-2.5ex]
%%XML element&Descripton\\
%%\hline
%%\end{tabular*}
%%\caption {\cmd{EASEMENTS\_17B} Table}
%%\label{table:easements_17b}
%%\end{table}

%% A simple table example
%% The htb attribute attempts to inline the table or figure
%% where you put it in the document
\begin{table}[htb]
\begin{center}
\begin{tabular*}{\textwidth}{@{}p{.25\textwidth}@{}p{.75\textwidth}}
\hline
\hline\\[-2.5ex]
XML element&Descripton\\
\hline
\hline\\[-1.5ex]   %% Trick to add whitespace after horizontal line
FNUM&US Survey file number\\ 
MERIDIAN&BLM meridian code\\
&\hspace{10pt}12 = Copper River\\
&\hspace{10pt}13 = Fairbanks\\
&\hspace{10pt}28 = Seward\\
&\hspace{10pt}44 = Kateel\\
&\hspace{10pt}45 = Umiat\\
TOWNSHIP&Five character Township code\\
RANGE&Five character Range code\\
PAGE&Survey page number 1,2,3,...\\
FILENAME&Relative path to file in direcory\\[1.5ex]
\hline
\end{tabular*}
\caption {USS XML index elements}
\label{table:uss-index}
\end{center}
\end{table}

%%
%% NOTE: tabular* takes table width argument
%%       Use of 'p' for left aligned column.
%%       Use of 'cp' to center column.
%%       The @{} works to remove an unwanted space.
%%
%% FIXME: Set width of centered columns?
%%
%% Links
%% =====
%% Getting to Grips with Latex - Tables
%% by Andrew Roberts
%% Source for @{\extracolsep{\fill} syntax.
%% http://www.andy-roberts.net/misc/latex/latextutorial4.html
%%
\begin{table}[htb]
\begin{center}
\begin{tabular*}{0.90\textwidth}{@{\extracolsep{\fill}}@{}l@{}c@{}c@{}c}
  \hline
  \hline\\[-2.5ex]
  col 1 & col 2 & col 3 & col 4 \\
  \hline
  \hline
  \hline\\[-1.5ex]
  item 1 & item 2 & item 3 & item 4 \\
  \hline
  item 1  & item 2  & item 3  & item 4  \\
  \hline
\end{tabular*}
\caption {Demo}
\label{table:demo-index}
\end{center}
\end{table}


%% -*- Mode: LaTeX -*-
%%
%% question.tex
%% Created Fri May 31 11:03:53 AKDT 2019
%% Copyright (C) 2019 by Raymond E. Marcil <marcilr@gmail.com>
%%
%% Questions
%%

%% ======================== Questions ============================
%% ======================== Questions ============================
\section{Questions}
\begin{itemize*}
     \item{} Break table column over two rows?
\end{itemize*}

%%\FIXME{None yet...}


%% -*- Mode: LaTeX -*-
%%
%% footnotes.tex
%% Created Fri May 31 12:44:40 AKDT 2019
%% Copyright (C) 2019 by Raymond E. Marcil <marcilr@gmail.com>
%%

%% ====================== Footnotes ======================
%% ====================== Footnotes ======================
\clearpage
\newpage
\subsection{Footnotes}
See my
footnote\footnote{\href{http://www.google.com/search?q=latex+footnotes}{Search}
google for footnotes.}
generated with:

\begin{verbatim}
  \footnote{\href{http://www.google.com/search?q=latex+footnotes}
  {Search google for footnotes.}}
\end{verbatim}

\noindent GoogleGuide --- Linking to Search
Results.\footnote{GoogleGuide --- \href{http://www.googleguide.com/linking.html}{Linking to Search Results.}}


%% -*- Mode: LaTeX -*-
%%
%%
%% endnotes.tex
%% Created Fri May 31 12:59:56 AKDT 2019
%% Copyright (C) 2019 by Raymond E. Marcil <marcilr@gmail.com>
%%

%% ======================== Endnotes =============================
%% ======================== Endnotes =============================
\clearpage
\newpage
\setcounter{secnumdepth}{0}
\section{Endnotes}

\begin{enumerate}

%% ====================== LaTeX Verbatim =========================
\item \LaTeX\ verbatim\\
\href{http://www.kfunigraz.ac.at/~binder/texhelp/ltx-79.html}
{http://www.kfunigraz.ac.at/~binder/texhelp/ltx-79.html}

\end{enumerate}

%% -*- Mode: LaTeX -*-
%%
%% appendix.tex
%% Created Fri May 31 13:03:23 AKDT 2019
%% Copyright (C) 2019 by Raymond E. Marcil <marcilr@gmail.com>
%%

%% ======================== Appendix =============================
%% ======================== Appendix =============================
%%
%% This will add a standard non-numbered Appendix to the document.
%% The next section Appendix chnages secnumdepth such that the Appendix
%% is not numbered but still displayed in the table of contents (TOC).
%%
%% Adding unnumbered sections to TOC
%% http://tex.stackexchange.com/questions/11668/adding-unnumbered-sections-to-toc
%% 
%% \section*{Appendix}
%% Need content here.
%%
\setcounter{secnumdepth}{0}
\section{Appendix}




%%
%% Main body of document
%%


%%\newpage
%%\lipsum[1]


\end{document}

%% Local Variables:
%% fill-column: 78
%% mode: auto-fill
%% compile-command: "make"
%% End:
